% Ad Atticum 11.13 --- headnote
% ad-atticum-11-13, mid-March 47 BC, Brundisium

\textit{Cicero to Atticus, written from Brundisium about
the middle of March 47 BC (the manuscript dateline:
\textit{Scr.\ Brundisi circ.\ med.\ m.\ Mart.\ a.\ 707
(47)}; month precision). Publius Siser has at last
delivered an answer from Atticus; Murena's freedman, who
was expected to bring one, has so far brought nothing.
Cicero corrects a piece of false news (that Quintus has
come to Syria) and asks his usual question about the
disposition of the men reaching Brundisium from
overseas --- so far, he reports, none of those who have
arrived has shown himself hostile, though he knows what
weight Atticus will put on the qualification. The
keynote of the section is a Stoic-tinged
self-recrimination: \textit{sola utilia mihi esse
videantur quae semper nolui} --- ``only those things
appear useful to me which I have always refused to want.''
The closing news of Publius Lentulus at Rhodes, his son at
Alexandria, and Gaius Cassius's move from Rhodes to
Alexandria are positional reports on the diaspora of
the Pompeian leadership.}

\textit{Section 2 is the heaviest. Brother Quintus has
sent an apology that is sharper in tone than his
earlier accusations: he is sorry to have offended
Atticus, but maintains that he was within his rights,
and sets out the grounds for his hatred ``in the
foulest fashion.'' Cicero reads the whole performance
as evidence that Quintus has only opened up his hatred
because he sees Cicero broken on every side. The
following sections handle estate business: the
co-heirs in the Fufidian inheritance ask only what is
fair and Atticus is to settle it as he sees best; about
the redemption of the Frusinate farm, Cicero is still
of the mind he was when matters stood better, but
asks Atticus to think about where the cash will come
from, since whatever liquid resources he had were
given over to Pompey when that still looked like
prudent loyalty --- a costly point of evidence against
brother Quintus's complaint that Cicero gave him
nothing, when Quintus neither asked nor saw the money.
The letter ends, as Book 11's letters increasingly do,
on the bare admission that grief keeps him from
writing more.}
